% !TEX root = ../TCC - CD e IA.tex
\chapter{Resultados}\label{Resultados}

% ==========================================================================
% ROTEIRO GERAL DO CAPÍTULO 6
% Propósito: este é o capítulo ACIONÁVEL. O cap. 5 apresentou a modelagem
% (o "motor"); aqui os resultados são organizados no funil COBRA-RACE e
% traduzidos em recomendações de growth. Fio condutor: sair da descrição
% ("o que os dados mostram") para a decisão ("o que a assessoria deve fazer").
% Não repetir a modelagem do cap. 5 — REFERENCIAR (\ref{}) e interpretar.
% ==========================================================================

% ROTEIRO (parágrafo de abertura do capítulo):
% - 1 parágrafo curto. Retomar a pergunta de pesquisa (do cap. 1) e anunciar
%   que este capítulo a responde aplicando o framework COBRA-RACE aos
%   resultados já obtidos no cap. 5.
% - Anunciar a estrutura do capítulo em 1 frase (as seções abaixo).

\section{Do Diagnóstico à Ação: Operacionalização do Funil de Engajamento}

% ROTEIRO:
% - Apresentar o framework COBRA-RACE já "montado" com OS SEUS dados.
% - Inserir uma TABELA (montar) que cruza: Estágio RACE | Nível COBRA |
%   métrica dos seus dados | cluster/figura correspondente.
%   Linhas: Reach->views(PC1); Act/Consumir->curtidas/cluster Padrão;
%   Convert/Contribuir->comentários+sentimento positivo;
%   Engage/Criar->comentários em debate (cluster Viral \ref{fig:sentiment-cluster-m1}).
% - Explicar que cada seção seguinte percorre um nível do funil.
% - Citar \cite{FEROZ2024} e \cite{MUNTINGA2011} ao apresentar a tabela.

\section{Segmentação de Publicações por Desempenho (Clusterização)}

% ROTEIRO:
% - Interpretar, à luz do growth, os 3 clusters de reels já obtidos (\ref{fig:clusters_pca}).
% - Retomar os números do cap.5/README: Cluster 0 "Padrão" (792 reels);
%   Cluster -1 "Viral" (12 reels, ~4.843 comentários, ~90k curtidas, 1,12M views);
%   Cluster 1 "Longo" (6 reels, ~721s de duração média).
% - NÃO reexplicar PCA/DBSCAN (está no cap.5) — referenciar \ref{tab:pca_loadings}.
% - Fechar cada cluster com sua LEITURA de growth (Padrão=aprovado,
%   Viral=debatido/polarizado, Longo=ignorado/indiferença).

\subsection{Clusterização de Publicações do Feed}

% ROTEIRO (PENDENTE DE EXECUÇÃO — você disse que vai gerar):
% - Rodar AutoClusterHPO também sobre os posts do FEED (hoje só há de reels).
% - Preencher com: nº de clusters encontrados, algoritmo eleito, score CVI,
%   perfil de cada cluster. Inserir figura análoga à \ref{fig:clusters_pca}.
% - Comparar com os clusters de reels: os padrões se repetem no feed?

\subsection{Clusterização de Perfis dos Governadores}

% ROTEIRO (PENDENTE DE EXECUÇÃO — você disse que vai gerar):
% - Rodar clusterização sobre as MÉTRICAS AGREGADAS por perfil
%   (tabela governor_engagement / \ref{tab:matriz_correlacao}).
% - Preencher: quantos grupos de governadores emergem, o que caracteriza
%   cada grupo (ex: "alto alcance/baixa conversão"), quais governadores em cada.
% - Este é o nível de segmentação à la Flesch \cite{FLESCH2017}: nomear os
%   grupos em termos funcionais se possível (loyalist-heavy, critic-heavy...).

\section{Qualificação dos Segmentos pelo Sentimento}

% ROTEIRO:
% - Cruzar clusters x sentimento (a etapa que "prova a tese", README seção 2.5).
% - Inserir a TABELA de distribuição de sentimento por cluster (montar a partir
%   das figuras \ref{fig:sentiment-cluster-m1}, \ref{fig:sentiment-cluster-0},
%   \ref{fig:sentiment-cluster-1}):
%   Cluster 0: 72,3% pos / 14,1% neu / 13,7% neg
%   Cluster -1 (Viral): 53,5% pos / 25,7% neu / 20,8% neg  <- polarizado
%   Cluster 1 (Longo): 60,4% pos / 27,1% neu / 12,5% neg   <- indiferença
% - Interpretar: alto engajamento (Viral) NÃO é alta aprovação. Este é o
%   achado central — conectá-lo explicitamente à crítica das vanity metrics
%   \cite{FEROZ2024} feita no referencial (seção 3.2).

\section{Discurso da Assessoria versus Reação do Público (Tópicos)}

% ROTEIRO:
% - Este é o contraste NOVO e mais rico. Precisa dos tópicos das legendas/
%   transcrições (discurso oficial) vs. tópicos dos comentários (público).
% - Comentários: retomar os 50 tópicos do BERTopic (\ref{fig:intertopic},
%   \ref{fig:hierarchy}, \ref{fig:heatmap}) já obtidos.
% - Legendas/transcrições (PENDENTE — vêm da API): rodar BERTopic sobre elas.
% - Montar a análise de GAP: temas que a assessoria pauta mas o público
%   ignora; temas que o público levanta mas a assessoria não aborda.
% - Este gap é a matéria-prima das recomendações da seção seguinte.

\section{Métricas de Crescimento e Priorização}

\subsection{North Star Metric: Engajamento Positivo Qualificado}

% ROTEIRO:
% - Definir a fórmula da NSM (a detalhar em bloco próprio — ver conversa).
%   Ideia: (comentários positivos / comentários totais) ponderado por alcance.
% - Calcular a NSM por perfil a partir de governor_sentiment + governor_engagement.
% - Contrastar o ranking por NSM com o ranking por engajamento bruto —
%   mostrar que mudam de ordem (prova de que a métrica importa).
% - Ancorar em \cite{FEROZ2024} e no precedente de Loyalty Rate (a citar).

\subsection{Retenção e Crescimento da Audiência (CMGR)}

% ROTEIRO:
% - Usar governor_engagement_history (tabela append, já existe) para calcular
%   crescimento ao longo do tempo, à la Flesch \cite{FLESCH2017}.
% - Se houver poucos pontos temporais ainda, declarar honestamente a limitação
%   e apresentar o método pronto para quando houver mais coletas.

\subsection{Priorização de Temas (Score ICE)}

% ROTEIRO:
% - Rankear os 50 tópicos por um score simples: Impacto (alcance do tópico x
%   sentimento positivo) x Confiança x Facilidade.
% - Entregar uma tabela "top temas a produzir" — saída diretamente acionável.

\section{Estudo de Caso Aplicado: Recomendações por Perfil}

% ROTEIRO:
% - AQUI o framework "paga" tudo. Escolher 2-3 governadores dos EXTREMOS já
%   identificados no cap.5:
%   * Clécio Luís (92,01% positivo, \ref{fig:sentiment-top5-positivo})
%   * Ibaneis (51,89% negativo, \ref{fig:sentiment-top5-negativo})
%   * (opcional) um caso mediano de contraste.
% - Para CADA um, aplicar o funil completo: em que estágio a audiência trava?
%   Qual o gap discurso x reação? O que a NSM revela? E então:
%   3-5 recomendações CONCRETAS de growth ("produzir mais X", "abandonar
%   formato Y", "o tema Z converte e está subexplorado").
% - Este é o núcleo do valor prático do TCC. Ser específico e acionável.

\section{Validação: Abordagem Híbrida versus Análises Isoladas}

% ROTEIRO (responde ao objetivo específico nº 6):
% - Argumentar, com os resultados em mãos, o que CADA técnica isolada NÃO
%   veria e que a integração revelou:
%   * só sentimento: veria "Ibaneis tem muita crítica" mas não em qual formato/tema.
%   * só tópicos: veria os temas mas não a valência nem o desempenho.
%   * só clusterização: veria grupos de performance mas não o conteúdo/reação.
% - Fechar com o achado que só o cruzamento permite: "engajamento != aprovação".